\documentclass[11pt]{article}

\usepackage[T1]{fontenc}
\usepackage{lmodern}
\usepackage{microtype}
\usepackage[margin=1in]{geometry}
\usepackage{amsmath,amssymb,amsthm,mathtools}
\usepackage{booktabs,tabularx,array}
\usepackage{enumitem}
\usepackage[style=authoryear,natbib=true,maxcitenames=2,maxbibnames=30,doi=true,url=true,isbn=false]{biblatex}
\addbibresource{davies_perturbation_companion_note.bib}
\usepackage{xcolor}
\usepackage[colorlinks=true,linkcolor=blue!55!black,citecolor=blue!55!black,urlcolor=blue!55!black]{hyperref}
\usepackage[nameinlink,noabbrev,capitalize]{cleveref}

\allowdisplaybreaks
\setlength{\parindent}{1.5em}
\setlength{\parskip}{0pt}
\setlength{\emergencystretch}{2em}

\newtheorem{theorem}{Theorem}[section]
\newtheorem{proposition}[theorem]{Proposition}
\newtheorem{lemma}[theorem]{Lemma}
\newtheorem{corollary}[theorem]{Corollary}
\newtheorem{remark}[theorem]{Remark}
\newtheorem{example}[theorem]{Example}
\crefname{theorem}{theorem}{theorems}
\Crefname{theorem}{Theorem}{Theorems}
\crefname{proposition}{proposition}{propositions}
\Crefname{proposition}{Proposition}{Propositions}
\crefname{lemma}{lemma}{lemmas}
\Crefname{lemma}{Lemma}{Lemmas}
\crefname{corollary}{corollary}{corollaries}
\Crefname{corollary}{Corollary}{Corollaries}
\crefname{remark}{remark}{remarks}
\Crefname{remark}{Remark}{Remarks}
\crefname{example}{example}{examples}
\Crefname{example}{Example}{Examples}

\newcommand{\R}{\mathbb R}
\newcommand{\E}{\mathbb E}
\newcommand{\cE}{\mathcal E}
\newcommand{\ip}[2]{\left\langle #1,#2\right\rangle}
\newcommand{\norm}[1]{\left\lVert #1\right\rVert}
\newcommand{\abs}[1]{\left\lvert #1\right\rvert}
\newcommand{\1}{\mathbf 1}
\newcommand{\D}{\mathsf D}

\title{A Detailed Davies-Perturbation Proof for Stationary Langevin Increments}
\author{Qian Qin}
\date{August 2026}

\begin{document}
\maketitle

\begin{abstract}
Let \((X_t)_{t\geq0}\) be stationary overdamped Langevin diffusion with
invariant distribution \(\pi(dx)\propto e^{-U(x)}dx\) and noise coefficient
\(\sqrt2\).  This note proves, for every \(v\in\R^d\), \(\lambda\in\R\),
and \(t\geq0\),
\[
 \E\exp\{\lambda\ip{v}{X_t-X_0}\}
 \leq \exp\{\lambda^2t\norm v^2\}.
\]
The result is a standard corollary of the Davies exponential-perturbation
estimate for symmetric semigroups.  The purpose of the note is to supply a
complete proof in the weighted Euclidean setting used by the companion MALA
paper and to explain the functional-analytic terms that a brief citation
normally suppresses: strongly continuous semigroup, generator domain,
core, semigroup equation, conjugated generator, and closed Dirichlet form.
A second part explains carefully when a deterministic-time process identity
may be evaluated at a stopping time, and applies those principles to the
localized stochastic-exponential calculation used in the MALA rejection
proof.
\end{abstract}

\tableofcontents

\section{The result and its place in the literature}\label{sec:intro}

Let \(U:\R^d\to\R\) be continuously differentiable, assume that
\(\nabla U\) is globally Lipschitz, and suppose
\[
 Z=\int_{\R^d}e^{-U(x)}\,dx<\infty.
\]
Define
\[
 \pi(dx)=Z^{-1}e^{-U(x)}\,dx.
\]
On a filtered probability space carrying a \(d\)-dimensional Brownian
motion \(B\), consider
\begin{equation}\label{eq:langevin}
 dX_t=-\nabla U(X_t)\,dt+\sqrt2\,dB_t,
 \qquad X_0\sim\pi.
\end{equation}
Global Lipschitz continuity gives a unique nonexplosive strong solution; see
\citet[Chapter~5]{KaratzasShreve1991}.  The process is stationary and
reversible with respect to \(\pi\).  One standard construction uses the
closed symmetric Dirichlet form
\begin{equation}\label{eq:dirichlet-form-intro}
 \cE(f,g)=\int_{\R^d}\ip{\nabla f(x)}{\nabla g(x)}\,\pi(dx),
\end{equation}
and identifies its diffusion with the solution of \eqref{eq:langevin}; see
\citet[Chapters~1 and~7]{FukushimaOshimaTakeda2011} and
\citet{BakryGentilLedoux2014}.

The estimate needed in the MALA paper is the following.

\begin{theorem}[Stationary linear-increment bound]\label{thm:increment}
For every \(v\in\R^d\), \(\lambda\in\R\), and \(t\geq0\),
\begin{equation}\label{eq:increment-mgf}
 \E\exp\{\lambda\ip{v}{X_t-X_0}\}
 \leq
 \exp\{\lambda^2t\norm v^2\}.
\end{equation}
\end{theorem}

The exact display \eqref{eq:increment-mgf} is not usually isolated as a
named theorem.  It is an immediate specialization of the exponential
perturbation, or Davies perturbation, method for symmetric semigroups.
Davies develops this method as part of the proof of Gaussian heat-kernel
bounds; see \citet[Chapter~3]{Davies1989}.  A systematic treatment through
closed forms and semigroups is given by
\citet[Sections~1.3--1.5 and~6.3]{Ouhabaz2005}.  The proof below records the
specialization with all constants and domains made explicit.

There is also a probabilistic proof.  The Lyons--Zheng decomposition writes
a stationary increment of a symmetric Markov process as one half of a
forward martingale minus one half of a reverse-time martingale.  Applied to
the coordinate function \(x\mapsto\ip{v}{x}\), it gives
\eqref{eq:increment-mgf} by Cauchy--Schwarz and Gaussian moment-generating
functions.  See \citet[Section~5.7, especially Theorem~5.7.1]{FukushimaOshimaTakeda2011},
the original paper of \citet{LyonsZheng1988}, and the later formulation of
\citet{KuwaeUemura1997}.  That route is developed in the first continuous-time
companion note \citep{QinContinuousTimeCompanion2026}.  The present note
concentrates on the shorter analytic route.

\section{A dictionary for semigroups and generators}\label{sec:dictionary}

This section gives the minimum functional analysis needed for the proof.
All Hilbert spaces below are real; this avoids irrelevant complex
conjugations.

\subsection{Strongly continuous semigroups}

Let \(H\) be a Hilbert space.  A family of bounded linear operators
\((T_t)_{t\geq0}\) on \(H\) is a \emph{strongly continuous semigroup}, or a
\(C_0\)-semigroup, if
\[
 T_0=I,\qquad T_{s+t}=T_sT_t,
 \qquad \lim_{t\downarrow0}\norm{T_tf-f}_H=0
 \quad\text{for every }f\in H.
\]
The last condition is continuity in the Hilbert-space norm, not pointwise
continuity in the state variable.

For a Markov process with invariant law \(\pi\), the transition operators
\[
 S_tf(x)=\E[f(X_t)\mid X_0=x]
\]
form a contraction semigroup on \(L^2(\pi)\).  In the reversible case they
are self-adjoint:
\[
 \ip{f}{S_tg}_{L^2(\pi)}
 =\ip{S_tf}{g}_{L^2(\pi)}.
\]
For the Langevin diffusion \eqref{eq:langevin}, strong continuity and
self-adjointness follow from the standard Dirichlet-form construction cited
above.

\subsection{The generator and its domain}

The generator \(A\) of a \(C_0\)-semigroup \((T_t)\) is defined by
\[
 Af=\lim_{h\downarrow0}\frac{T_hf-f}{h},
\]
but this limit need not exist for every \(f\in H\).  Its domain is
\begin{equation}\label{eq:generator-domain}
 \D(A)=
 \left\{
 f\in H:
 \lim_{h\downarrow0}\frac{T_hf-f}{h}
 \text{ exists in }H
 \right\}.
\end{equation}
The generator is generally unbounded, but it is closed and its domain is
dense in \(H\).

The basic semigroup differentiation theorem is the following.  It is a
standard result in any text on \(C_0\)-semigroups; see
\citet[Chapter~1]{Pazy1983} or
\citet[Chapter~II]{EngelNagel2000}.

\begin{proposition}[Semigroup equation]\label{prop:semigroup-equation}
If \(f\in\D(A)\), then \(T_tf\in\D(A)\) for every \(t\geq0\), and the map
\(t\mapsto T_tf\) is differentiable in \(H\), with
\begin{equation}\label{eq:semigroup-equation}
 \frac{d}{dt}T_tf=AT_tf=T_tAf.
\end{equation}
\end{proposition}

Thus the ``semigroup equation'' is the infinite-dimensional analogue of
\[
 \frac{d}{dt}e^{tA}f=Ae^{tA}f
\]
for a matrix \(A\).  The domain condition is what makes the derivative
legitimate.

\subsection{What a core is, and why the proof below does not need one}

A subspace \(\mathcal C\subseteq\D(A)\) is a \emph{core} for \(A\) if it is
dense in \(\D(A)\) for the graph norm
\[
 \norm f_A:=\norm f_H+\norm{Af}_H.
\]
Equivalently, the closure of the restricted operator
\(A|_{\mathcal C}\) is the full generator \(A\).  A core is useful when one
first identifies a differential expression on a convenient class, such as
\(C_c^\infty(\R^d)\), and then wishes to conclude that its closure is the
actual generator.

The proof in \cref{sec:davies} does not say ``take a function in a core.''
Instead, it works directly with the full generator domain
\(\D(A_\psi)\) of the conjugated semigroup.  The semigroup equation is then
available from \cref{prop:semigroup-equation}.  The only differential
calculation is performed through the already closed Dirichlet form, where
bounded Lipschitz multiplication is legitimate.

\section{The Langevin Dirichlet form}\label{sec:form}

Let \(H=L^2(\pi)\).  The closure of
\eqref{eq:dirichlet-form-intro}, initially defined on
\(C_c^\infty(\R^d)\), has domain
\begin{equation}\label{eq:form-domain}
 \D(\cE)=
 \left\{
 f\in L^2(\pi):
 f\text{ has a weak gradient and }
 \int\norm{\nabla f}^2\,d\pi<\infty
 \right\}.
\end{equation}
This is the weighted Sobolev space \(H^1(\pi)\).  The associated
self-adjoint generator is denoted by \(\mathcal L\).  Our sign convention is
that \(\mathcal L\) is nonpositive.  For smooth compactly supported
functions,
\[
 \mathcal Lf=\Delta f-\ip{\nabla U}{\nabla f}.
\]
The generator and form are related by
\begin{equation}\label{eq:form-generator-relation}
 \ip{g}{\mathcal Lf}_{L^2(\pi)}=-\cE(g,f),
 \qquad f\in\D(\mathcal L),\quad g\in\D(\cE).
\end{equation}
In particular, \(\D(\mathcal L)\subseteq\D(\cE)\).

We will repeatedly multiply form-domain functions by bounded Lipschitz
functions.  The needed product rule is standard for weak derivatives, but
we record it explicitly.

\begin{lemma}[Bounded Lipschitz multipliers]\label{lem:multipliers}
Let \(a:\R^d\to\R\) be bounded and belong to \(W^{1,\infty}(\R^d)\).  If
\(f\in\D(\cE)\), then \(af\in\D(\cE)\) and
\begin{equation}\label{eq:weak-product}
 \nabla(af)=a\nabla f+f\nabla a
\end{equation}
in the weak sense.  In particular, if \(\psi\in W^{1,\infty}\) is bounded,
then multiplication by \(e^{\psi}\) and by \(e^{-\psi}\) maps
\(\D(\cE)\) into itself.
\end{lemma}

\begin{proof}
Choose smooth functions \(f_n\) converging to \(f\) in the norm
\[
 \norm f_{\D(\cE)}^2=\norm f_{L^2(\pi)}^2+\cE(f,f).
\]
The weak product rule for a Sobolev function and a bounded Lipschitz
function gives \eqref{eq:weak-product}; alternatively, approximate \(a\) by
smooth Lipschitz functions with the same uniform bounds.  The estimates
\[
 \norm{af}_{L^2(\pi)}\leq\norm a_\infty\norm f_{L^2(\pi)},
\]
and
\[
 \norm{\nabla(af)}_{L^2(\pi)}
 \leq
 \norm a_\infty\norm{\nabla f}_{L^2(\pi)}
 +\norm{\nabla a}_\infty\norm f_{L^2(\pi)}
\]
show that \(af\in\D(\cE)\).  Taking \(a=e^{\pm\psi}\) gives the last
assertion because \(\psi\) is bounded and Lipschitz.  Standard Sobolev
multiplier facts of this type are discussed, for example, in
\citet[Chapter~1]{Ouhabaz2005}.
\end{proof}

\section{The Davies perturbation estimate}\label{sec:davies}

Let \(\psi\in W^{1,\infty}(\R^d)\) be bounded.  Define multiplication
operators
\[
 M_{\psi}f=e^{\psi}f,
 \qquad
 M_{-\psi}f=e^{-\psi}f.
\]
They are bounded inverses on \(L^2(\pi)\).  Define the conjugated operators
\begin{equation}\label{eq:conjugated-semigroup}
 T_t^\psi=M_{-\psi}S_tM_\psi,
 \qquad
 T_t^\psi f=e^{-\psi}S_t(e^\psi f).
\end{equation}

\begin{lemma}[The conjugated generator]\label{lem:conjugated-generator}
The family \((T_t^\psi)_{t\geq0}\) is a strongly continuous semigroup.  Its
generator \(\mathcal L_\psi\) has domain
\begin{equation}\label{eq:conjugated-domain}
 \D(\mathcal L_\psi)
 =M_{-\psi}\D(\mathcal L)
 =\{f:e^\psi f\in\D(\mathcal L)\},
\end{equation}
and
\begin{equation}\label{eq:conjugated-generator}
 \mathcal L_\psi f=e^{-\psi}\mathcal L(e^\psi f).
\end{equation}
\end{lemma}

\begin{proof}
The semigroup property follows by cancellation of the middle multiplication
operators:
\[
 T_s^\psi T_t^\psi
 =M_{-\psi}S_sM_\psi M_{-\psi}S_tM_\psi
 =M_{-\psi}S_{s+t}M_\psi
 =T_{s+t}^\psi.
\]
Strong continuity follows from boundedness of the multiplication operators:
\[
 \norm{T_t^\psi f-f}_2
 \leq
 \norm{e^{-\psi}}_\infty
 \norm{S_t(e^\psi f)-e^\psi f}_2
 \longrightarrow0.
\]
Finally,
\[
 \frac{T_t^\psi f-f}{t}
 =e^{-\psi}
 \frac{S_t(e^\psi f)-e^\psi f}{t}.
\]
Because multiplication by \(e^{-\psi}\) is a bounded invertible operator,
the limit on the left exists in \(L^2(\pi)\) if and only if the limit inside
on the right exists.  By the definition of the generator domain, this is
equivalent to \(e^\psi f\in\D(\mathcal L)\), and the limit is
\eqref{eq:conjugated-generator}.
\end{proof}

\begin{lemma}[Numerical-range estimate]\label{lem:numerical-range}
For every \(f\in\D(\mathcal L_\psi)\),
\begin{equation}\label{eq:numerical-range}
 \ip{f}{\mathcal L_\psi f}_{L^2(\pi)}
 \leq
 \norm{\nabla\psi}_\infty^2\norm f_{L^2(\pi)}^2.
\end{equation}
\end{lemma}

\begin{proof}
Let \(f\in\D(\mathcal L_\psi)\) and set \(g=e^\psi f\).  Then
\(g\in\D(\mathcal L)\subseteq\D(\cE)\).  By
\cref{lem:multipliers}, both \(f=e^{-\psi}g\) and
\(e^{-\psi}f=e^{-2\psi}g\) belong to \(\D(\cE)\).  Therefore the
form-generator identity \eqref{eq:form-generator-relation} applies:
\begin{align*}
 \ip{f}{\mathcal L_\psi f}_{L^2(\pi)}
 &=\ip{e^{-\psi}f}{\mathcal L(e^\psi f)}_{L^2(\pi)}\\
 &=-\cE(e^{-\psi}f,e^\psi f).
\end{align*}
Using the weak product rule,
\[
 \nabla(e^{-\psi}f)=e^{-\psi}(\nabla f-f\nabla\psi),
 \qquad
 \nabla(e^\psi f)=e^\psi(\nabla f+f\nabla\psi).
\]
The cross terms cancel pointwise, so
\begin{align*}
 \cE(e^{-\psi}f,e^\psi f)
 &=\int
 \ip{\nabla f-f\nabla\psi}{\nabla f+f\nabla\psi}
 \,d\pi\\
 &=\int\norm{\nabla f}^2\,d\pi
   -\int f^2\norm{\nabla\psi}^2\,d\pi.
\end{align*}
Consequently,
\[
 \ip{f}{\mathcal L_\psi f}
 =-\int\norm{\nabla f}^2\,d\pi
  +\int f^2\norm{\nabla\psi}^2\,d\pi
 \leq
 \norm{\nabla\psi}_\infty^2\norm f_2^2,
\]
which is \eqref{eq:numerical-range}.
\end{proof}

\begin{theorem}[Davies perturbation estimate]\label{thm:davies}
For every bounded \(\psi\in W^{1,\infty}(\R^d)\) and every \(t\geq0\),
\begin{equation}\label{eq:davies}
 \norm{e^{-\psi}S_t(e^\psi\,\cdot)}_{L^2(\pi)\to L^2(\pi)}
 \leq
 \exp\{t\norm{\nabla\psi}_\infty^2\}.
\end{equation}
\end{theorem}

\begin{proof}
First take \(f\in\D(\mathcal L_\psi)\) and put
\(u(t)=T_t^\psi f\).  By \cref{prop:semigroup-equation},
\(u(t)\in\D(\mathcal L_\psi)\) and
\[
 u'(t)=\mathcal L_\psi u(t)
\]
in \(L^2(\pi)\).  The Hilbert-space chain rule therefore gives
\begin{align*}
 \frac{d}{dt}\norm{u(t)}_2^2
 &=2\ip{u(t)}{u'(t)}_2\\
 &=2\ip{u(t)}{\mathcal L_\psi u(t)}_2\\
 &\leq2\norm{\nabla\psi}_\infty^2\norm{u(t)}_2^2,
\end{align*}
where the last line is \cref{lem:numerical-range}.  Gronwall's inequality
gives
\begin{equation}\label{eq:davies-domain}
 \norm{T_t^\psi f}_2
 \leq
 e^{t\norm{\nabla\psi}_\infty^2}\norm f_2,
 \qquad f\in\D(\mathcal L_\psi).
\end{equation}
The domain \(\D(\mathcal L)\) is dense in \(L^2(\pi)\), and multiplication
by \(e^{-\psi}\) is bounded and invertible.  Hence
\(\D(\mathcal L_\psi)=e^{-\psi}\D(\mathcal L)\) is also dense.  For an
arbitrary \(f\in L^2(\pi)\), choose
\(f_n\in\D(\mathcal L_\psi)\) with \(f_n\to f\) in \(L^2(\pi)\).
Since \(T_t^\psi\) is bounded for each fixed \(t\),
\(T_t^\psi f_n\to T_t^\psi f\).  Passing to the limit in
\eqref{eq:davies-domain} proves \eqref{eq:davies}.
\end{proof}

\begin{remark}[Where the usual core argument went]\label{rem:core}
A shorter presentation often verifies
\cref{lem:numerical-range} first on \(C_c^\infty\), says that this class is
a core, differentiates \(T_t^\psi f\), and then extends by closure.  That is
valid once the core statement and the passage to the conjugated generator
are justified.  The proof above avoids that compressed language: the
Dirichlet form supplies the identity on the full generator domain, and
\cref{lem:conjugated-generator} identifies the conjugated domain directly.
\end{remark}

\section{Proof of the stationary increment bound}\label{sec:application}

We now apply \cref{thm:davies}.  Choose smooth functions
\(\chi_R:\R\to\R\) such that
\[
 \chi_R(s)=s\quad\text{when }\abs s\leq R,
 \qquad
 \abs{\chi_R'(s)}\leq1,
 \qquad
 \chi_R\text{ is bounded}.
\]
For example, one may begin with a smooth function whose derivative is one
on \([-1,1]\), vanishes outside \([-2,2]\), and lies between zero and one,
and then rescale.  Set
\[
 \psi_R(x)=\lambda\chi_R(\ip{v}{x}).
\]
Then \(\psi_R\) is bounded and
\[
 \norm{\nabla\psi_R}_\infty\leq\abs\lambda\norm v.
\]
Stationarity, the Markov property, and the definition of the semigroup give
\begin{align}
 \E e^{\psi_R(X_t)-\psi_R(X_0)}
 &=\int e^{-\psi_R(x)}S_t(e^{\psi_R})(x)\,\pi(dx)\notag\\
 &=\ip{1}{T_t^{\psi_R}1}_{L^2(\pi)}.\label{eq:stationary-conjugation}
\end{align}
Because \(\pi\) is a probability measure,
\(\norm1_{L^2(\pi)}=1\).  Cauchy--Schwarz and \cref{thm:davies} imply
\begin{align*}
 \E e^{\psi_R(X_t)-\psi_R(X_0)}
 &\leq\norm{T_t^{\psi_R}}_{2\to2}\\
 &\leq e^{\lambda^2t\norm v^2}.
\end{align*}
For every sample path,
\[
 \psi_R(X_t)-\psi_R(X_0)
 \longrightarrow
 \lambda\ip{v}{X_t-X_0}
 \qquad\text{as }R\to\infty.
\]
The exponential is nonnegative.  Fatou's lemma therefore gives
\[
 \E e^{\lambda\ip{v}{X_t-X_0}}
 \leq e^{\lambda^2t\norm v^2},
\]
which proves \cref{thm:increment}.

\begin{remark}[Interpretation of the constant]
The right-hand side is the moment-generating function of a centered Gaussian
with variance \(2t\norm v^2\).  Thus \(X_t-X_0\) is sub-Gaussian with
covariance proxy \(2tI_d\).  The conclusion does not assert that the
increment itself is Gaussian.  The constant reflects the normalization
\(\sqrt2\,dB_t\) in \eqref{eq:langevin}; with diffusion coefficient
\(dB_t\), the corresponding exponent would be divided by two.
\end{remark}

\section{A concise Lyons--Zheng cross-check}\label{sec:lz}

For completeness, we record why the probabilistic literature gives the same
estimate.  Let \(f_v(x)=\ip{v}{x}\).  The coordinate function belongs to the
Dirichlet-form domain whenever it is square integrable under \(\pi\), and
\[
 \cE(f_v,f_v)=\norm v^2.
\]
For a symmetric process started from its symmetrizing measure, the
Lyons--Zheng decomposition at the terminal time \(t\) gives
\[
 f_v(X_t)-f_v(X_0)
 =\frac12\left(M_t^v-\widehat M_t^{v,t}\right),
\]
where \(M^v\) is the forward martingale part and
\(\widehat M^{v,t}\) is the reverse-time martingale part.  In the present
Langevin diffusion,
\[
 [M^v]_t=[\widehat M^{v,t}]_t=2t\norm v^2,
\]
and each terminal martingale is centered Gaussian because its integrand is
the deterministic vector \(\sqrt2v\) in its respective Brownian filtration.
They are generally dependent.  Cauchy--Schwarz, rather than independence,
gives
\begin{align*}
 \E e^{\lambda(f_v(X_t)-f_v(X_0))}
 &\leq
 \left(\E e^{\lambda M_t^v}\right)^{1/2}
 \left(\E e^{-\lambda\widehat M_t^{v,t}}\right)^{1/2}\\
 &=e^{\lambda^2t\norm v^2}.
\end{align*}
This reproduces \cref{thm:increment}.  The abstract theorem and the passage
to unbounded coordinate functions are discussed in the references cited in
\cref{sec:intro}.

\section{Stopping times: what may and may not be substituted}\label{sec:stopping}

The MALA appendix repeatedly writes a process at \(t\wedge\tau\), where
\(\tau\) is a stopping time.  This is legitimate, but not for the reason
that an identity holds ``for every deterministic \(t\), almost surely'' if
the exceptional null set is allowed to depend on \(t\).  This section makes
the distinction precise.  Standard references for the stopping results
used here are \citet[Chapter~3]{KaratzasShreve1991},
\citet[Chapter~IV]{RevuzYor1999}, and
\citet[Chapter~II]{Protter2005}.

\subsection{Why fixed-time almost-sure equality is not enough}

Suppose that, for every deterministic \(t\), random variables \(Y_t\) and
\(Z_t\) satisfy \(Y_t=Z_t\) almost surely.  This statement alone does not
imply \(Y_\tau=Z_\tau\) for a random time \(\tau\), because the exceptional
null set may depend on \(t\).  A simple illustration is obtained on
\(\Omega=[0,1]\) with Lebesgue measure: set
\[
 Y_t(\omega)=\1_{\{\omega=t\}},\qquad Z_t(\omega)=0,
 \qquad \tau(\omega)=\omega.
\]
For each fixed \(t\), \(Y_t=Z_t\) almost surely, but
\(Y_\tau=1\) everywhere.  If one equips this space with the completed
filtration generated by the events \({\tau\leq s}\) up to time \(t\), then
\(\tau\) is a stopping time (and \(Y_t\) is adapted), so the pathology is
not removed merely by the stopping-time property.  This example is
deliberately irregular, but it shows why the order of the quantifiers
matters.

What is sufficient is an \emph{indistinguishable process identity}: there
is one event \(\Omega_0\) of probability one such that
\[
 Y_t(\omega)=Z_t(\omega)
 \quad\text{for every }t\geq0\text{ and every }\omega\in\Omega_0.
\]
Then the equality may plainly be evaluated at \(t=\tau(\omega)\).  If two
continuous processes agree almost surely at every rational time, continuity
and countability imply that they are indistinguishable.  This is why
continuous versions are especially convenient.

\subsection{Standard stopped-process identities}

Let \(M\) be a continuous local martingale and \(\tau\) a stopping time.
The stopped process is
\[
 M_t^\tau=M_{t\wedge\tau}.
\]
The following are standard process identities, valid up to
indistinguishability:
\begin{equation}\label{eq:stopped-bracket}
 [M^\tau]_t=[M]_{t\wedge\tau},
\end{equation}
and, if
\(M_t=\int_0^t\ip{H_s}{dB_s}\),
\begin{equation}\label{eq:stopped-integral}
 M_{t\wedge\tau}
 =\int_0^t\1_{\{s\leq\tau\}}\ip{H_s}{dB_s}.
\end{equation}
The distinction between \(s<\tau\) and \(s\leq\tau\) is immaterial for a
continuous It\^o integral.  Formula \eqref{eq:stopped-integral} can first be
checked for simple predictable integrands and then extended by the It\^o
isometry.  Formula \eqref{eq:stopped-bracket} follows from the defining
quadratic-variation theorem or from \eqref{eq:stopped-integral} when the
martingale is a Brownian integral.

Ordinary algebraic identities between continuous processes may also be
stopped.  For example, if
\[
 D_t=\exp\{M_t-V_t/2\}
\]
for every \(t\) on one probability-one event, then
\[
 D_{t\wedge\tau}
 =\exp\{M_{t\wedge\tau}-V_{t\wedge\tau}/2\}
\]
on the same event.  If \(V\) is increasing, then
\[
 V_{t\wedge\tau}\leq V_t
\]
pathwise.  Neither statement uses optional sampling.

By contrast, an expectation identity such as \(\E M_t=\E M_0\) for every
deterministic \(t\) does not automatically imply
\(\E M_\tau=\E M_0\).  That conclusion is an optional-sampling statement
and requires its own hypotheses.  The localized MALA proof uses the
process identities above, not an unjustified optional-sampling step.

\subsection{Line-by-line application to the path-likelihood lemma}

For clarity, we match the stopping operations in the MALA proof to the
results above.  Let
\[
 M_t=\int_0^t\ip{\theta_s}{dB_s},
 \qquad
 V_t=[M]_t=\int_0^t\norm{\theta_s}^2\,ds,
 \qquad
 D_t=e^{M_t-V_t/2},
\]
and let
\[
 \tau_n=\inf\{t:V_t\geq n\}.
\]

\begin{enumerate}[label=\textnormal{(\roman*)},leftmargin=2.7em]
\item The identity
\[
 [M^{\tau_n}]_t=V_{t\wedge\tau_n}
\]
is exactly \eqref{eq:stopped-bracket}.  It is not obtained by taking a
fixed-time equality and informally replacing \(t\) by a random time.

\item The factorization of \(D_{h\wedge\tau_n}^q\) into a stopped
stochastic exponential and an ordinary exponential is pointwise algebra
applied to the indistinguishable process identity defining \(D\).  The
inequality
\(V_{h\wedge\tau_n}\leq V_h\) is pathwise because \(V\) is increasing.

\item It\^o's formula gives the process identity
\[
 D_t-1=\int_0^tD_s\ip{\theta_s}{dB_s},
 \qquad t\geq0,
\]
up to indistinguishability.  Stopping this identity at a further stopping
time \(\sigma_n\) and using \eqref{eq:stopped-integral} gives
\[
 D_{t\wedge\sigma_n}-1
 =\int_0^t\1_{\{s\leq\sigma_n\}}
 D_s\ip{\theta_s}{dB_s}.
\]
This is the square-integrable martingale to which BDG is applied.

\item The stopping is finally removed pathwise.  On a fixed compact time
interval, the continuous processes \(V\) and \(D\) are bounded on every
sample path.  Hence the stopping times eventually exceed the terminal
horizon, and the stopped variables equal the unstopped variables for all
large \(n\).  Fatou's lemma then transfers the uniform moment bound to the
unstopped limit.
\end{enumerate}

Thus the answer to the motivating question is: \emph{partly, but the
relevant assertion is stronger}.  One needs process identities that hold
simultaneously for all times outside one null set, or the standard theorem
identifying the stopped process, stochastic integral, and quadratic
variation.  Merely knowing a separate almost-sure statement for each fixed
\(t\) would not suffice.

\section{Which parts are imported and which are specialized}\label{sec:status}

The following table separates literature results from calculations specific
to the MALA project.

\begin{center}
\begin{tabularx}{\textwidth}{>{\raggedright\arraybackslash}p{0.24\textwidth}
>{\raggedright\arraybackslash}X
>{\raggedright\arraybackslash}X}
\toprule
Ingredient & Imported result & Specialization in the MALA paper \\
\midrule
Stationary linear increments
& Davies' exponential perturbation estimate; alternatively, the
Lyons--Zheng decomposition
& Apply the estimate to a truncated linear function and retain the exact
constant for the \(\sqrt2\)-Langevin normalization. \\
\addlinespace
Quadratic and integrated increments
& Gaussian randomization is a standard device for turning linear
sub-Gaussian bounds into quadratic-form bounds; compare
\citet[Theorem~1 and Remark~2]{HsuKakadeZhang2012}
& Use an auxiliary standard Gaussian, then Jensen's inequality in time, to
obtain the exact moment-generating and \(L^p\) bounds for
\(\int_0^h\norm{X_t-X_0}^2dt\). \\
\addlinespace
Path likelihood
& Novikov, Girsanov, stopped stochastic exponentials, Doob, and
quantitative BDG are standard.  General reverse-H\"older theory for
continuous stochastic exponentials is developed by
\citet{Kazamaki1994}; related MALA diffusion comparisons appear in
\citet[Appendix~A]{ChewiEtAl2021}
& Choose the drift difference, prove the dimension- and moment-dependent
step restriction, and exploit
\(V_h\leq L^2\int_0^h\norm{X_t-X_0}^2dt/2\). \\
\addlinespace
Endpoint conditioning
& The likelihood ratio of a statistic is the conditional expectation of
the full likelihood ratio; see, for example, \citet{Kallenberg2021}
& Condition the path likelihood on \((X_0,X_h)\) and identify the exact and
Euler endpoint measures. \\
\addlinespace
Metropolis rejection
& The accepted Metropolis--Hastings flow is the meet of proposal flow and
its transpose \citep{Tierney1998,BilleraDiaconis2001}
& Use the symmetric exact endpoint measure as a dominating reference and
bound the pointwise rejection function. \\
\bottomrule
\end{tabularx}
\end{center}

The first row is the only lemma in the MALA appendix that is essentially a
direct specialization of one standard literature theorem.  The remaining
lemmas combine standard tools in a way adapted to the rejection problem;
they should cite those tools and nearby diffusion-comparison literature,
but should not be described as verbatim quotations.

\printbibliography

\end{document}
